\section{Joint squarefreeness beyond the actual root interval}
\label{sqw-section}

The following result proves simultaneous first-depth membership on a
growing interval of primes for a majority of actual roots. It does not
assume that supported primes are usually simple. The larger-prime tail
and the exceptional roots remain separate problems.

Let $n\ge5$ be prime, put $B=n^4$, and for integers $B\le k<2B$ write
\[
 a=3k,\qquad w_k=a+\zeta,\qquad Q_k=a^2+a+1,\qquad
 T_n(k)=|P(w_k^n)|,\qquad t_k=\log c_n(k).
\]
These are primitive unramified roots, and $T_1(k)=a(a+1)$.
At a supported prime $q\ge5$, the rank is
$d_q=\operatorname{ord}((w_k/\bar w_k)^3\bmod q)$.

\begin{theorem}[An actual joint repeated-prime bound]\label{sqw-joint}
For a real number $U>B$, let $\mathcal B(n,U)$ be the set of indices
in the block for which some prime $B<q\le U$ satisfies $q^2\mid T_n(k)$.
Then
\begin{equation}\label{sqw-bad-bound}
 \frac{|\mathcal B(n,U)|}{B}
 \le \frac{3n}{B}+
       \frac{6U}{B\log(U/(3n))}.
\end{equation}
In particular, with
\[
 U_n=\left\lfloor\frac{B\log n}{\log\log n}\right\rfloor,
\]
one has $U_n>B$ for all sufficiently large prime $n$, and
$|\mathcal B(n,U_n)|/B=O(1/\log\log n)$.
Outside this single exceptional set, every supported prime in
$(B,U_n]$ has exact depth one simultaneously.
\end{theorem}
\begin{proof}
A supported prime does not divide $Q_k$. Since $q>B>n$, the
homogeneous valuation law has no exponent-lifting term and its rank
divides $n$. Thus its rank is one or $n$.

Rank one cannot produce a repeated prime in this interval. Its depth
equals its depth in $T_1=a(a+1)$. These two factors are coprime, so
$q^2\mid T_1$ would force $q^2$ to divide one individual factor.
Each is less than $6B+1$, whereas $q^2>B^2>6B+1$ for $B\ge625$.
The comparison is with the individual factors, not their product.

Every bad prime therefore has rank $n$. The reviewed finite-field
rank calculation gives at most $3\varphi(n)$ simple roots in the
$k$-variable modulo $q$. Each has a unique lift modulo $q^2$,
with rank unchanged. Counting actual integers in the block gives
\[
 \#\{k:d_q=n,\ q^2\mid T_n(k)\}
 \le3\varphi(n)(B/q^2+1).
\]
Every eligible prime lies in one of the two classes $\pm1\bmod 3n$.
The first term in the union bound, after division by $B$, is at most
\[
 3\varphi(n)\sum_{m>B}m^{-2}\le3n/B.
\]
The two-progression Brun--Titchmarsh bound used in the earlier
rank-window proof gives
\[
 \#\{q\le U:q\equiv\pm1\bmod 3n\}
 \le\frac{4U}{\varphi(3n)\log(U/(3n))}.
\]
Here $U>B>3n$ and $\varphi(3n)=2\varphi(n)$, yielding the second
term of \eqref{sqw-bad-bound}. No independence of primes or uniform
distribution of actual roots has been assumed. Finally
$U_n/B\sim\log n/\log\log n$ and
$\log(U_n/(3n))\sim3\log n$, giving the stated asymptotic bound.
The exceptional set is defined by the existence of any repeated
prime, so its complement gives the simultaneous assertion.
\end{proof}

The exact rank and simple-lifting inputs are the ordinary results
underlying the sixth-round actual-root window. The sieve input is
the same real-endpoint specialization of Montgomery--Vaughan's
Theorem~2, equation~(1.10), already verified in the seventh-round
sieved-window proof. No stronger version of that input is invoked.

\begin{lemma}[Actual negative credit in the window]\label{sqw-credit}
For $k\notin\mathcal B(n,U_n)$, define, over supported primes,
\[
 M_{\rm win}=\sum_{B<q\le U_n}v_q(T_n)\log q,\qquad
 W_{\rm far}=\sum_{q>U_n}(v_q(T_n)-3)\log q,\qquad
 H_{\rm low}=\sum_{q\le B}(v_q(T_n)-3)\log q.
\]
Then
\begin{align}
 M_{\rm win}&=\sum_{\substack{B<q\le U_n\\q\mid T_n}}\log q,
 \label{sqw-window-rad}\\
 \mathcal J(T_n)&=W_{\rm far}-2M_{\rm win}+H_{\rm low}
 \le W_{\rm far}-2M_{\rm win}+L_B(T_n).
 \label{sqw-credit-identity}
\end{align}
\end{lemma}
\begin{proof}
The preceding theorem makes each supported depth in the window equal
to one. Its signed contribution is therefore $-2\log q$. Partition
the global supported-prime sum at $B$ and $U_n$, and bound the
remaining small-prime signed contribution by its full mass.
\end{proof}

This is proved credit for actual roots, but there is no positive lower
bound for $M_{\rm win}$. The window may contain no supported primes.
The next result also shows that its typical mass is sublinear.

\begin{theorem}[The remaining packet still carries the full mass]
\label{sqw-far-mass}
Put $\epsilon_n=(\log\log n)^{-1/2}$. There is an absolute constant
$C$ and, for every sufficiently large prime $n$, a set $\mathcal H_n$
of actual indices of size at least $(1-C\epsilon_n)B$, such that every
member has the joint squarefree property above and
\begin{align}
 L_{U_n}(T_n(k))/t_k&\le\epsilon_n,&
 0\le M_{\rm win}(k)/t_k&\le\epsilon_n, \label{sqw-small-mass}\\
 3-\frac4n-\epsilon_n
 &\le\frac1{t_k}\sum_{q>U_n}v_q(T_n(k))\log q
 \le3. \label{sqw-top-mass}
\end{align}
Every supported prime in \eqref{sqw-top-mass} has exact rank $n$,
and its valuation is its first depth.
\end{theorem}
\begin{proof}
Apply the full-mass mean in Section~\ref{fml-section} at $U_n$.
For prime $n$, $\tau(n)=2$ and $\sigma(n)/n=1+1/n$.
The three terms of that bound are respectively
$O(\log n/n)$, $O(1/\log\log n)$ and $O(1/n)$.
The elementary estimate of Section~\ref{ebi-section} adds at most
$20/n$ for primes two and three. Hence
\[
 \frac1B\sum_{B\le k<2B}\frac{L_{U_n}(T_n(k))}{t_k}
 \le\frac{C}{\log\log n}.
\]
Markov at $\epsilon_n$ excludes at most $C\epsilon_n$ of the
relative mass of indices. Intersect with the complement of
$\mathcal B(n,U_n)$, whose exceptional fraction is
$O(\epsilon_n^2)$. This proves \eqref{sqw-small-mass}.
The elementary angular estimate
$3-4/n\le\log T_n/t_k\le3$ gives \eqref{sqw-top-mass}
by subtraction.

Eventually $U_n>6B+1$, so no prime above $U_n$ divides either old
input factor $a,a+1$. Rank one is therefore absent altogether.
Also $q>n$, so there is no exponent lifting. The homogeneous
valuation law identifies its depth with its rank-$n$ first depth.
\end{proof}

For an actual sequence in $\mathcal H_n$, the condition
\[
 (W_{\rm far}-2M_{\rm win})_+/t_k\longrightarrow0
\]
is sufficient for $\limsup\mathcal J(T_n)/t_k\le0$ and hence for the
restricted ABC estimate, for each fixed positive tolerance at
sufficiently large indices. This is a conditional conclusion.
On the proved set, the available window credit is itself $o(t_k)$,
whereas the farther packet carries $(3-o(1))t_k$ of full valuation
mass. The squarefree window does not pay for a positive linear
signed cost in that packet. Neither that signed cost nor the
excluded roots have been controlled. The analytic membership
theorem here is an independently reviewed ordinary proof; no
completed Lean proof of its sieve or exceptional-set statements
is asserted.
